\SourceLeaf{349}%
the curved surface is specified by an equation relating \(x,y,z\). It becomes more complicated still when the surface is specified by expressing \(x,y,z\) as functions of two new variables \(p,q\). In this last case the expression contains fifteen elements, namely the first and second partial derivatives of \(x,y,z\) with respect to \(p\) and \(q\). Its chief importance is that it leads to another expression, which ranks among the most remarkable results of this theory. When the curved surface is specified in this way, the general expression for any line element on it,
\[
\text{or for}\quad\BareRoot(\mathrm{d}x^2+\mathrm{d}y^2+\mathrm{d}z^2),\quad\text{takes the form}\quad\BareRoot(E\mathrm{d}x^2+2F\mathrm{d}x\mathord{.}\mathrm{d}y+G\mathrm{d}y^2)
\]
where \(E,F.\ G\) are in turn functions of \(p\) and \(q\). The new expression for the measure of curvature contains only these quantities and their first and second partial derivatives. Thus the general expression for a line element alone suffices to determine the measure of curvature, independently of the expressions for the coordinates \(x,y,z\) themselves. A direct consequence is the remarkable theorem that if a curved surface, or a patch of it, can be developed onto another surface, the measure of curvature at every point remains unchanged after the development. A special case follows. On a curved surface that can be developed onto a plane, the measure of curvature is everywhere \(=0\). This immediately gives the characteristic equation of surfaces that can be developed onto a plane, where \(z\) is regarded as a function of \(x\) and \(y\),
\[
\frac{\mathrm{dd}z}{\mathrm{d}x^2}\mathord{.}\frac{\mathrm{dd}z}{\mathrm{d}y^2}-(\frac{\mathrm{dd}z}{\mathrm{d}x\mathord{.}\mathrm{d}y})^2=0
\]
an equation known for a long time, although in the author's judgment a proof with the required rigor had yet to be given.

These results suggest a new way of looking at the theory of curved surfaces, opening up a wide field of inquiry that is still wholly uncultivated. If one treats surfaces as bodies with one vanishing dimension, setting aside their role as boundaries of bodies, and regards them as flexible and inextensible, one sees that two essentially different kinds of
